\documentclass[12pt]{article}
\usepackage{amsmath}
\usepackage{amssymb}
\usepackage{physics}
\usepackage{hyperref}

\newcommand{\ds}{\displaystyle}
\newcommand{\ip}[2]{\langle #1 \mid #2 \rangle}

\title{The Relativistic Maxwell Equations}
\author{}
\date{}

\begin{document}
\maketitle


% ============================================================
\section{Introduction}
% ============================================================

In this segment, we will explore the derivation of the Maxwell equations' transformation for
a general velocity boost in any given direction.

Traditionally, the derivation of relativistic Maxwell equations is tackled through
differential geometry. Einstein's seminal work utilized standard multivariate calculus, albeit
limited to boosts aligned with coordinate axes. Extending this to boosts in arbitrary
directions significantly complicates the calculus involved, a topic seldom addressed in
existing literature or online resources.

Our journey through the relativistic Maxwell equations will be rooted in multivariate
calculus. The foundational definitions and indispensable formulas will be laid out in the
opening section, setting the stage for the intricate derivations that follow.

The subsequent section delves into the transformation of key differential operators---divergence
and curl---under the Lorentz transformation, a pivotal step in understanding the equations'
behaviour under relativistic conditions.

Culminating in the final section, we will tackle the derivation of the relativistic Maxwell
equations within a realm devoid of charges and currents, showcasing the equations' purest
form in a charge-free vacuum.

\subsection{Definitions and Useful Formulas from Multivariate Calculus}

We will consider a stationary system of reference where the coordinates of an event are
represented by the four-vector $x_1, x_2, x_3, t$, and a moving frame where the same event
has coordinates $x'_1, x'_2, x'_3, t'$. The fourth coordinate in both four-vectors has units
of time, not space, as the mathematics of the Maxwell equations derivation becomes a little
less notation-heavy this way.

A vector field $\vb{A}$ is a three-dimensional valued function acting on a four-vector,
\begin{equation}
  \begin{gathered}
    \vb{A}(x_1, x_2, x_3, t) \\
    \vb{A}: \mathbb{R}^4 \to \mathbb{R}^3
  \end{gathered}
\end{equation}

A few differential operators of $\vb{A}$ are extensively used in the derivation below and
deserve to be explicitly defined to minimize confusion when the notation will get a little
heavy, as it often does in the derivation of the relativistic Maxwell equations.

\begin{itemize}
  \item The Jacobian on the spatial dimensions and time:
  \begin{equation}
    \nabla \vb{A}_{x,t} =
    \begin{pmatrix}
      \pdv{A_1}{x_1} & \pdv{A_1}{x_2} & \pdv{A_1}{x_3} & \pdv{A_1}{t} \\[6pt]
      \pdv{A_2}{x_1} & \pdv{A_2}{x_2} & \pdv{A_2}{x_3} & \pdv{A_2}{t} \\[6pt]
      \pdv{A_3}{x_1} & \pdv{A_3}{x_2} & \pdv{A_3}{x_3} & \pdv{A_3}{t}
    \end{pmatrix}
  \end{equation}

  \item The Jacobian on the spatial dimensions:
  \begin{equation}
    \nabla \vb{A}_{x} =
    \begin{pmatrix}
      \pdv{A_1}{x_1} & \pdv{A_1}{x_2} & \pdv{A_1}{x_3} \\[6pt]
      \pdv{A_2}{x_1} & \pdv{A_2}{x_2} & \pdv{A_2}{x_3} \\[6pt]
      \pdv{A_3}{x_1} & \pdv{A_3}{x_2} & \pdv{A_3}{x_3}
    \end{pmatrix}
  \end{equation}

  \item The Jacobian for the spatial dimension applied to a vector, typically a velocity
        $\vb{v}$:
  \begin{equation}
    \nabla \vb{A}_x \cdot \vb{v} =
    \begin{pmatrix}
      \pdv{A_1}{x_1}v_1 & \pdv{A_1}{x_2}v_2 & \pdv{A_1}{x_3}v_3 \\[6pt]
      \pdv{A_2}{x_1}v_1 & \pdv{A_2}{x_2}v_2 & \pdv{A_2}{x_3}v_3 \\[6pt]
      \pdv{A_3}{x_1}v_1 & \pdv{A_3}{x_2}v_2 & \pdv{A_3}{x_3}v_3
    \end{pmatrix}
  \end{equation}

  \item The divergence of a vector field $\vb{A}$:
  \begin{equation}
    \div \vb{A} = \pdv{A_1}{x_1} + \pdv{A_2}{x_2} + \pdv{A_3}{x_3}
  \end{equation}

  \item And finally the curl of a vector field $\vb{A}$:
  \begin{equation}
    \curl \vb{A} =
      \left(\pdv{A_3}{x_2} - \pdv{A_2}{x_3}\right)\vb{e}_1
      - \left(\pdv{A_3}{x_1} - \pdv{A_1}{x_3}\right)\vb{e}_2
      + \left(\pdv{A_2}{x_1} - \pdv{A_1}{x_2}\right)\vb{e}_3
  \end{equation}
  where $\vb{e}_1$, $\vb{e}_2$ and $\vb{e}_3$ are unit vectors along the three spatial axes.
\end{itemize}

A few identities from calculus are key to derive the relativistic Maxwell equations, and they
are listed below:

\begin{itemize}
  \item The divergence and the curl of the cross product of a constant vector $\vb{v}$ and a
        field $\vb{A}$:
  \begin{equation}
    \begin{aligned}
      \curl(\vb{v} \cross \vb{A}) &= -\nabla_x \vb{A} \cdot \vb{v} + (\div \vb{A})\,\vb{v} \\
      \div(\vb{v} \cross \vb{A})  &= -\ip{\vb{v}}{\curl \vb{A}}
    \end{aligned}
    \label{eq:divergence_curl_cross_product}
  \end{equation}

  \item The curl of the triple product of a constant vector $\vb{v}$ with a field $\vb{A}$,
        which has two different expressions, the latter derived with the use of the equations
        above:
  \begin{equation}
    \begin{aligned}
      \curl\bigl(\vb{v} \cross (\vb{v} \cross \vb{A})\bigr)
        &= \curl\bigl(\ip{\vb{v}}{\vb{A}}\vb{v} - v^2\vb{A}\bigr) \\[4pt]
      \curl\bigl(\vb{v} \cross (\vb{v} \cross \vb{A})\bigr)
        &= -\nabla_x(\vb{v} \cross \vb{A}) \cdot \vb{v}
           + \div(\vb{v} \cross \vb{A})\,\vb{v} \\
        &= -\vb{v} \cross (\nabla_x \vb{A} \cdot \vb{v})
           - \ip{\vb{v}}{\curl \vb{A}}\,\vb{v}
    \end{aligned}
    \label{eq:curl_triple_product}
  \end{equation}
\end{itemize}

% ============================================================
\section{Transformation of Differential Operators}
% ============================================================

Before working on the Maxwell equations, we need to understand how the differential operators
transform between $F$ and $F'$. We consider a field $\vb{A}$ in $F$,
\begin{equation}
  \begin{gathered}
    \vb{A} = A(x_1, x_2, x_3, t) \\
    \vb{A}: \mathbb{R}^4 \to \mathbb{R}^3
  \end{gathered}
\end{equation}
which we compose with the inverse of the Lorentz transformation from $F'$ to $F$ (see
\textit{The Transformation of Time and Space}) obtaining a field $\vb{A}'$ defined on the
space $F'$,
\begin{align}
  \vb{A}' &= \vb{A} \circ \boldsymbol{\Lambda}(-\vb{v}) \\
  \vb{A}' &= \vb{A}'(\vb{x}', t') = \vb{A}\bigl(\boldsymbol{\Lambda}(-\vb{v})(\vb{x}', t')\bigr) \nonumber \\
  \vb{A}' &: \mathbb{R}^4 \to \mathbb{R}^3 \nonumber
\end{align}

Since the Maxwell equations are in $F$ coordinates, we need to understand how the partial
derivatives of the field $\vb{A}'$ transform under the Lorentz transformation, which we are
able to compute from the Jacobian chain rule,
\begin{equation}
  \nabla_{x',t'}\vb{A}' = \nabla_{x,t}\vb{A} \cdot \nabla_{x',t'}\bigl(\boldsymbol{\Lambda}(-\vb{v})(\vb{x}',t')\bigr)
\end{equation}
where $\nabla_{x',t'}\vb{A}'$ and $\nabla_{x,t}\vb{A}$ are the Jacobians of $\vb{A}'$ and
$\vb{A}$ respectively, computed in their own coordinate systems.

The Lorentz transformation is linear, hence the Jacobian is $\Lambda$, and the chain rule
becomes explicitly in coordinate form:
\begin{equation}
  \begin{pmatrix}
    \ds\pdv{A'_1}{x'_1} & \ds\pdv{A'_1}{x'_2} & \ds\pdv{A'_1}{x'_3} & \ds\pdv{A'_1}{t'} \\[8pt]
    \ds\pdv{A'_2}{x'_1} & \ds\pdv{A'_2}{x'_2} & \ds\pdv{A'_2}{x'_3} & \ds\pdv{A'_2}{t'} \\[8pt]
    \ds\pdv{A'_3}{x'_1} & \ds\pdv{A'_3}{x'_2} & \ds\pdv{A'_3}{x'_3} & \ds\pdv{A'_3}{t'}
  \end{pmatrix}
  =
  \begin{pmatrix}
    \ds\pdv{A_1}{x_1} & \ds\pdv{A_1}{x_2} & \ds\pdv{A_1}{x_3} & \ds\pdv{A_1}{t} \\[8pt]
    \ds\pdv{A_2}{x_1} & \ds\pdv{A_2}{x_2} & \ds\pdv{A_2}{x_3} & \ds\pdv{A_2}{t} \\[8pt]
    \ds\pdv{A_3}{x_1} & \ds\pdv{A_3}{x_2} & \ds\pdv{A_3}{x_3} & \ds\pdv{A_3}{t}
  \end{pmatrix}
  \cdot\, \Lambda(-\vb{v})
\end{equation}

To get to the partial derivatives of $\vb{A}$ we need to invert the equation above,
\begin{equation}
  \begin{pmatrix}
    \ds\pdv{A'_1}{x'_1} & \ds\pdv{A'_1}{x'_2} & \ds\pdv{A'_1}{x'_3} & \ds\pdv{A'_1}{t'} \\[8pt]
    \ds\pdv{A'_2}{x'_1} & \ds\pdv{A'_2}{x'_2} & \ds\pdv{A'_2}{x'_3} & \ds\pdv{A'_2}{t'} \\[8pt]
    \ds\pdv{A'_3}{x'_1} & \ds\pdv{A'_3}{x'_2} & \ds\pdv{A'_3}{x'_3} & \ds\pdv{A'_3}{t'}
  \end{pmatrix}
  \cdot\,\Lambda(\vb{v})
  =
  \begin{pmatrix}
    \ds\pdv{A_1}{x_1} & \ds\pdv{A_1}{x_2} & \ds\pdv{A_1}{x_3} & \ds\pdv{A_1}{t} \\[8pt]
    \ds\pdv{A_2}{x_1} & \ds\pdv{A_2}{x_2} & \ds\pdv{A_2}{x_3} & \ds\pdv{A_2}{t} \\[8pt]
    \ds\pdv{A_3}{x_1} & \ds\pdv{A_3}{x_2} & \ds\pdv{A_3}{x_3} & \ds\pdv{A_3}{t}
  \end{pmatrix}
\end{equation}

The bra-ket formalism helps us here once again to compute the product of the Jacobian
matrices. Any row $j$ of $\nabla_{x',t'}\vb{A}' \cdot \Lambda(\vb{v})$ is the vector
\begin{equation}
  \begin{aligned}
    &\bra{\ds\pdv{A'_j}{x'_1},\, \ds\pdv{A'_j}{x'_2},\, \ds\pdv{A'_j}{x'_3},\, \ds\pdv{A'_j}{t'}}
    \left[
      \frac{\beta}{v^2}\ket{\vb{v},0}\bra{\vb{v},0}
      - \beta\ket{\vb{v},0}\bra{\vb{0},1} \right.\\
    &\left.
      + \left(\vb{I} - \frac{1}{v^2}\ket{\vb{v},0}\bra{\vb{v},0}\right)
      - \frac{\beta}{c^2}\ket{\vb{0},1}\bra{\vb{v},0}
      + \beta\ket{\vb{0},1}\bra{\vb{0},1}
    \right]
  \end{aligned}
\end{equation}
where $\vb{I}$ is the identity operator, and reads,
\begin{equation}
  \begin{aligned}
    &\frac{\beta}{v^2}\ip{\nabla_{x'}\vb{A}'_j}{\vb{v}}\bra{\vb{v},0}
    - \beta\ip{\nabla_{x'}\vb{A}'_j}{\vb{v}}\bra{\vb{0},1}
    + \bra{\nabla_{x'}\vb{A}'_j,\,0} \\
    &- \frac{1}{v^2}\ip{\nabla_{x'}\vb{A}'_j}{\vb{v}}\bra{\vb{v},0}
    - \frac{\beta}{c^2}\ds\pdv{A'_j}{t'}\bra{\vb{v},0}
    + \beta\ds\pdv{A'_j}{t'}\bra{\vb{0},1}
  \end{aligned}
\end{equation}

The spatial gradient of any component of the field $\vb{A}$ computed at
$(\vb{x},t) = \Lambda(-\vb{v})(\vb{x}',t')$ is,
\begin{equation}
  \nabla_x \vb{A}_j = \nabla_{x'}\vb{A}'_j
    + \frac{\beta-1}{v^2}\ip{\nabla_{x'}\vb{A}'_j}{\vb{v}}\bra{\vb{v}}
    - \frac{\beta}{c^2}\ds\pdv{A'_j}{t'}\bra{\vb{v}}
  \label{eq:gradient_in_F}
\end{equation}
and the time derivative is
\begin{equation}
  \pdv{A_j}{t} = -\beta\ip{\nabla_{x'}\vb{A}'_j}{\vb{v}} + \beta\pdv{A'_j}{t'}
  \label{eq:time_derivative_in_F}
\end{equation}
where,
\begin{align}
  \nabla_{x'}\vb{A}'_j &= \bra{\ds\pdv{A'_j}{x'_1},\, \ds\pdv{A'_j}{x'_2},\, \ds\pdv{A'_j}{x'_3}} \\
  \nabla_x \vb{A}_j    &= \bra{\ds\pdv{A_j}{x_1},\, \ds\pdv{A_j}{x_2},\, \ds\pdv{A_j}{x_3}}
\end{align}

\subsection{Divergence in $F$ and $F'$}

We now use \eqref{eq:gradient_in_F} to compute the divergence of the field $\vb{A}$ in $F$,
\begin{align}
  \pdv{A_1}{x_1} &= \pdv{A'_1}{x'_1}
    + \frac{\beta-1}{v^2}\ip{\nabla_{x'}\vb{A}'_1}{\vb{v}}v_1
    - \frac{\beta}{c^2}\pdv{A'_1}{t'}v_1 \\
  \pdv{A_2}{x_2} &= \pdv{A'_2}{x'_2}
    + \frac{\beta-1}{v^2}\ip{\nabla_{x'}\vb{A}'_2}{\vb{v}}v_2
    - \frac{\beta}{c^2}\pdv{A'_2}{t'}v_2 \\
  \pdv{A_3}{x_3} &= \pdv{A'_3}{x'_3}
    + \frac{\beta-1}{v^2}\ip{\nabla_{x'}\vb{A}'_3}{\vb{v}}v_3
    - \frac{\beta}{c^2}\pdv{A'_3}{t'}v_3
\end{align}

To derive an expression for the field divergence, we expand the scalar products in the
equation above and group the terms on the partial derivative along each coordinate. We will
work out the details for the first equation, as the others are identically derived by just
changing the indexes.
\begin{equation}
  \begin{aligned}
    \pdv{A_1}{x_1} &= \pdv{A'_1}{x'_1}
      + \frac{\beta-1}{v^2}\ip{\nabla_{x'}\vb{A}'_1}{\vb{v}}v_1
      - \frac{\beta}{c^2}\pdv{A'_1}{t'}v_1 \\
    &= \pdv{A'_1}{x'_1}
      + \frac{\beta-1}{v^2}\!\left(
          \pdv{A'_1}{x'_1}v_1 + \pdv{A'_1}{x'_2}v_2 + \pdv{A'_1}{x'_3}v_3
        \right)v_1
      - \frac{\beta}{c^2}\pdv{A'_1}{t'}v_1 \\
    &= \left(1 + \frac{\beta-1}{v^2}v_1^2\right)\pdv{A'_1}{x'_1}
      + \frac{\beta-1}{v^2}\!\left(
          \pdv{(A'_1 v_1)}{x'_2}v_2 + \pdv{(A'_1 v_1)}{x'_3}v_3
        \right)
      - \frac{\beta}{c^2}\pdv{(A'_1 v_1)}{t'}
  \end{aligned}
\end{equation}

The divergence of the field in $F$ is,
\begin{equation}
  \begin{aligned}
    &\pdv{A_1}{x_1} + \pdv{A_2}{x_2} + \pdv{A_3}{x_3} \\
    &\quad= \left(1 + \frac{\beta-1}{v^2}v_1^2\right)\pdv{A'_1}{x'_1}
      + \frac{\beta-1}{v^2}\!\left(\pdv{(A'_1 v_1)}{x'_2}v_2 + \pdv{(A'_1 v_1)}{x'_3}v_3\right)
      - \frac{\beta}{c^2}\pdv{(A'_1 v_1)}{t'} \\
    &\quad+ \left(1 + \frac{\beta-1}{v^2}v_2^2\right)\pdv{A'_2}{x'_2}
      + \frac{\beta-1}{v^2}\!\left(\pdv{(A'_2 v_2)}{x'_1}v_1 + \pdv{(A'_2 v_2)}{x'_3}v_3\right)
      - \frac{\beta}{c^2}\pdv{(A'_2 v_2)}{t'} \\
    &\quad+ \left(1 + \frac{\beta-1}{v^2}v_3^2\right)\pdv{A'_3}{x'_3}
      + \frac{\beta-1}{v^2}\!\left(\pdv{(A'_3 v_3)}{x'_1}v_1 + \pdv{(A'_3 v_3)}{x'_2}v_2\right)
      - \frac{\beta}{c^2}\pdv{(A'_3 v_3)}{t'}
  \end{aligned}
\end{equation}

By grouping the terms by the components of the field $\vb{A}$, we obtain the divergence of
the field $\vb{A}$ in $F$,
\begin{equation}
  \begin{aligned}
    &\pdv{A_1}{x_1} + \pdv{A_2}{x_2} + \pdv{A_3}{x_3} \\
    &= \pdv{A'_1}{x'_1} + \pdv{A'_2}{x'_2} + \pdv{A'_3}{x'_3} \\
    &\quad + \frac{\beta-1}{v^2}\!\left(\pdv{A'_1}{x'_1}v_1 + \pdv{A'_1}{x'_2}v_2 + \pdv{A'_1}{x'_3}v_3\right)v_1 \\
    &\quad + \frac{\beta-1}{v^2}\!\left(\pdv{A'_2}{x'_1}v_1 + \pdv{A'_2}{x'_2}v_2 + \pdv{A'_2}{x'_3}v_3\right)v_2 \\
    &\quad + \frac{\beta-1}{v^2}\!\left(\pdv{A'_3}{x'_1}v_1 + \pdv{A'_3}{x'_2}v_2 + \pdv{A'_3}{x'_3}v_3\right)v_3 \\
    &\quad - \frac{\beta}{c^2}\pdv{(A'_1 v_1 + A'_2 v_2 + A'_3 v_3)}{t'}
  \end{aligned}
\end{equation}

We get the compact equation for the transformation of the divergence of $\vb{A}$ to $\vb{A}'$,
\begin{equation}
  \div \vb{A} = \div \vb{A}'
    + \frac{\beta-1}{v^2}\ip{\nabla_{x'}\vb{A}' \cdot \vb{v}}{\vb{v}}
    - \frac{\beta}{c^2}\pdv{\ip{\vb{A}'}{\vb{v}}}{t'}
\end{equation}

If we consider a divergence-free field $\vb{A}$ in $F$, the equation above gives an important
condition on the divergence of $\vb{A}'$,
\begin{equation}
  \div \vb{A}' = -\frac{\beta-1}{v^2}\ip{\nabla_{x'}\vb{A}' \cdot \vb{v}}{\vb{v}}
    + \frac{\beta}{c^2}\pdv{\ip{\vb{A}'}{\vb{v}}}{t'}
  \label{eq:divergence_in_Fprime}
\end{equation}

\subsection{Curl in $F$ and $F'$}

The other operator needed for the Maxwell equations is the curl,
\begin{equation}
  \begin{aligned}
    (\curl \vb{A})_1 &= \pdv{A_3}{x_2} - \pdv{A_2}{x_3} \\
      &= \pdv{A'_3}{x'_2}
        + \frac{\beta-1}{v^2}\ip{\nabla_{x'}\vb{A}'_3}{\vb{v}}v_2
        - \frac{\beta}{c^2}\pdv{A'_3}{t'}v_2 \\
      &\quad - \pdv{A'_2}{x'_3}
        - \frac{\beta-1}{v^2}\ip{\nabla_{x'}\vb{A}'_2}{\vb{v}}v_3
        + \frac{\beta}{c^2}\pdv{A'_2}{t'}v_3 \\[6pt]
    (\curl \vb{A})_2 &= -\pdv{A_3}{x_1} + \pdv{A_1}{x_3} \\
      &= -\pdv{A'_3}{x'_1}
        - \frac{\beta-1}{v^2}\ip{\nabla_{x'}\vb{A}'_3}{\vb{v}}v_1
        - \frac{\beta}{c^2}\pdv{A'_3}{t'}v_1 \\
      &\quad - \pdv{A'_1}{x'_3}
        - \frac{\beta-1}{v^2}\ip{\nabla_{x'}\vb{A}'_1}{\vb{v}}v_3
        + \frac{\beta}{c^2}\pdv{A'_1}{t'}v_3 \\[6pt]
    (\curl \vb{A})_3 &= \pdv{A_2}{x_1} - \pdv{A_1}{x_2} \\
      &= \pdv{A'_2}{x'_1}
        + \frac{\beta-1}{v^2}\ip{\nabla_{x'}\vb{A}'_2}{\vb{v}}v_1
        - \frac{\beta}{c^2}\pdv{A'_2}{t'}v_1 \\
      &\quad - \pdv{A'_1}{x'_2}
        - \frac{\beta-1}{v^2}\ip{\nabla_{x'}\vb{A}'_1}{\vb{v}}v_2
        + \frac{\beta}{c^2}\pdv{A'_1}{t'}v_2
  \end{aligned}
\end{equation}

which is in vector form easily recognisable as,
\begin{equation}
  \curl \vb{A} = \curl \vb{A}'
    + \frac{\beta-1}{v^2}\,\vb{v} \cross (\nabla_{x'}\vb{A}' \cdot \vb{v})
    - \frac{\beta}{c^2}\pdv{\vb{v} \cross \vb{A}'}{t'}
  \label{eq:curl_in_F}
\end{equation}

% ============================================================
\section{The Relativistic Maxwell Equations}
% ============================================================

We consider an electric field $\vb{E}$ and a magnetic field $\vb{B}$ in a region of the
system of reference $F$ free of charges and currents. The Maxwell equations in this case are
\begin{align}
  \pdv{\vb{B}}{t} &= -\curl \vb{E} \nonumber \\
  \pdv{\vb{E}}{t} &= c^2 \curl \vb{B} \nonumber \\
  \div \vb{B}     &= 0                \nonumber \\
  \div \vb{E}     &= 0
  \label{eq:maxwell_in_F}
\end{align}

The fields $\vb{B}'$ and $\vb{E}'$, obtained by applying the composition of the Lorentz
transformation are,
\begin{align}
  \vb{B}' &= \vb{B}'(\vb{x}',t') = \vb{B}\bigl(\Lambda(-\vb{v})(\vb{x}',t')\bigr) \\
  \vb{E}' &= \vb{E}'(\vb{x}',t') = \vb{E}\bigl(\Lambda(-\vb{v})(\vb{x}',t')\bigr)
\end{align}

and fulfill the equations obtained by applying directly the Lorentz transformation to
equations \eqref{eq:maxwell_in_F} (see section \textbf{Transformation of Differential
Operators}),
\begin{align}
  -\beta(\nabla_{x'}\vb{B}' \cdot \vb{v}) + \beta\pdv{\vb{B}'}{t'}
    &= -\curl \vb{E}'
       - \frac{\beta-1}{v^2}\,\vb{v} \cross (\nabla_{x'}\vb{E}' \cdot \vb{v})
       + \frac{\beta}{c^2}\pdv{\vb{v} \cross \vb{E}'}{t'} \label{eq:maxwell_transformed} \\
  -\frac{\beta}{c^2}(\nabla_{x'}\vb{E}' \cdot \vb{v}) + \frac{\beta}{c^2}\pdv{\vb{E}'}{t'}
    &= \curl \vb{B}'
       + \frac{\beta-1}{v^2}\,\vb{v} \cross (\nabla_{x'}\vb{B}' \cdot \vb{v})
       - \frac{\beta}{c^2}\pdv{\vb{v} \cross \vb{B}'}{t'} \nonumber
\end{align}

Since the equations at hand do not conform to the standard form of Maxwell's equations, the
fields denoted by $\vb{B}'$ and $\vb{E}'$ do not represent the transformed electric and
magnetic fields we seek. Substantial effort is required to manipulate these equations into the
correct form that yields the desired fields.

We outline below the steps required to derive the relativistic Maxwell equations in a vacuum,
starting with the equations above and manipulating them to obtain the correct form.

\subsection{Step 1: Divergence Free Fields}

The scalar product of \eqref{eq:maxwell_in_F} with the velocity $\vb{v}$ gives,
\begin{align}
  -\beta\ip{\nabla_{x'}\vb{B}' \cdot \vb{v}}{\vb{v}}
    + \beta\pdv{\ip{\vb{B}'}{\vb{v}}}{t'}
    &= -\ip{\curl \vb{E}'}{\vb{v}} \\
  -\frac{\beta}{c^2}\ip{\nabla_{x'}\vb{E}' \cdot \vb{v}}{\vb{v}}
    + \frac{\beta}{c^2}\pdv{\ip{\vb{E}'}{\vb{v}}}{t'}
    &= \ip{\curl \vb{B}'}{\vb{v}}
\end{align}

together with the divergence-free condition of the fields \eqref{eq:divergence_in_Fprime}, it
gives a very useful expression for the divergence of $\vb{B}'$ and $\vb{E}'$,
\begin{align}
  \div \vb{B}' &= \frac{\beta-1}{\beta v^2}
    \left(\pdv{\ip{\vb{B}'}{\vb{v}}}{t'} - \ip{\curl \vb{E}'}{\vb{v}}\right) \\
  \div \vb{E}' &= \frac{\beta-1}{\beta v^2}
    \left(\pdv{\ip{\vb{E}'}{\vb{v}}}{t'} + c^2\ip{\curl \vb{B}'}{\vb{v}}\right)
\end{align}

The terms $\nabla_{x'}\vb{B}' \cdot \vb{v}$ and $\nabla_{x'}\vb{E}' \cdot \vb{v}$ in the
left side of \eqref{eq:maxwell_transformed} can be substituted from
\eqref{eq:divergence_curl_cross_product}, and using the divergence of $\vb{B}'$ and $\vb{E}'$
just calculated, we obtain,
\begin{align}
  \pdv{}{t'}\!\left(\beta\,\vb{B}' - \frac{\beta-1}{v^2}\ip{\vb{B}'}{\vb{v}}\vb{v}
      - \frac{\beta}{c^2}\vb{v} \cross \vb{E}'\right)
    &= -\beta\curl(\vb{v} \cross \vb{B}') - \beta\curl \vb{E}' \nonumber \\
    &\quad - \frac{\beta-1}{v^2}\ip{\curl \vb{E}'}{\vb{v}}\vb{v}
      - \frac{\beta-1}{v^2}\,\vb{v} \cross (\nabla_{x'}\vb{E}' \cdot \vb{v}) \label{eq:divergence_free_fields} \\[6pt]
  \pdv{}{t'}\!\left(\beta\,\vb{E}' - \frac{\beta-1}{v^2}\ip{\vb{E}'}{\vb{v}}\vb{v}
      + \beta\,\vb{v} \cross \vb{B}'\right)
    &= -\beta\curl(\vb{v} \cross \vb{E}') + \curl(c^2\vb{B}') \nonumber \\
    &\quad + \frac{\beta-1}{v^2}\ip{\curl(c^2\vb{B}')}{\vb{v}}\vb{v}
      + \frac{\beta-1}{v^2}\,\vb{v} \cross (\nabla_{x'}(c^2\vb{B}') \cdot \vb{v}) \nonumber
\end{align}

We are nearing the solution: it is now necessary to apply the curl operation to the triple
product.

\subsection{Step 2: The Curl of the Triple Product}

The equations \eqref{eq:divergence_free_fields} are getting us closer to the solution. We
have on the left side the partial derivative with respect to the time $t'$ of two fields. On
the right side, we have a few curl operators, which is what we need, but we still have terms
that are not in curl form:
\begin{equation}
  -\frac{\beta-1}{v^2}\ip{\curl \vb{E}'}{\vb{v}}\vb{v}
  - \frac{\beta-1}{v^2}\,\vb{v} \cross (\nabla_{x'}\vb{E}' \cdot \vb{v})
\end{equation}

The key to the solution here is to use the curl of the triple product, which is given by
equation \eqref{eq:curl_triple_product},
\begin{equation}
  \curl\bigl(\vb{v} \cross (\vb{v} \cross \vb{E}')\bigr)
  = \curl\bigl(\ip{\vb{v}}{\vb{E}'}\vb{v} - v^2\vb{E}'\bigr)
  = -\vb{v} \cross (\nabla_x \vb{E}' \cdot \vb{v}) - \ip{\vb{v}}{\curl \vb{E}'}\vb{v}
\end{equation}

which substituted in \eqref{eq:divergence_free_fields} yields finally equations in the same
operator form as equations \eqref{eq:maxwell_in_F}, but now in the coordinates of $F'$:
\begin{align}
  \pdv{}{t'}\!\left(\beta\,\vb{B}' - (\beta-1)\ip{\vb{B}'}{\frac{\vb{v}}{v}}\frac{\vb{v}}{v}
      - \frac{\beta}{c^2}\vb{v} \cross \vb{E}'\right)
    &= -\curl\!\left(\beta\,\vb{E}'
        - (\beta-1)\ip{\vb{E}'}{\tfrac{\vb{v}}{v}}\tfrac{\vb{v}}{v}
        + \beta\,\vb{v} \cross \vb{B}'\right) \label{eq:pre_final_maxwell} \\[6pt]
  \pdv{}{t'}\!\left(\beta\,\vb{E}' - (\beta-1)\ip{\vb{E}'}{\frac{\vb{v}}{v}}\frac{\vb{v}}{v}
      + \beta\,\vb{v} \cross \vb{B}'\right)
    &= -c^2\curl\!\left(\beta\,\vb{B}'
        - (\beta-1)\ip{\vb{B}'}{\tfrac{\vb{v}}{v}}\tfrac{\vb{v}}{v}
        - \frac{\beta}{c^2}\vb{v} \cross \vb{E}'\right) \nonumber
\end{align}

Hence, by the principle of relativity, the electric and magnetic fields in $F'$ must be
\begin{align}
  \vb{B}_F &= \phi(\vb{v})\left(
      \beta\,\vb{B}' - (\beta-1)\ip{\vb{B}'}{\tfrac{\vb{v}}{v}}\tfrac{\vb{v}}{v}
      - \frac{\beta}{c^2}\vb{v} \cross \vb{E}'
    \right) \\
  \vb{E}_F &= \phi(\vb{v})\left(
      \beta\,\vb{E}' - (\beta-1)\ip{\vb{E}'}{\tfrac{\vb{v}}{v}}\tfrac{\vb{v}}{v}
      + \beta\,\vb{v} \cross \vb{B}'
    \right)
\end{align}
where $\phi(\vb{v})$ is a scalar function depending only on the velocity.

The condition on the possible values of $\phi(\vb{v})$ can be found by requiring that the
inversion of the electric and magnetic fields in $F'$ is the identity,
\begin{align}
  \vb{B} &= \phi(-\vb{v})\left(
      \beta\,\vb{B}_F - (\beta-1)\ip{\vb{B}_F}{\tfrac{\vb{v}}{v}}\tfrac{\vb{v}}{v}
      + \frac{\beta}{c^2}\vb{v} \cross \vb{E}_F
    \right) \\
  \vb{E} &= \phi(-\vb{v})\left(
      \beta\,\vb{E}_F - (\beta-1)\ip{\vb{E}_F}{\tfrac{\vb{v}}{v}}\tfrac{\vb{v}}{v}
      - \beta\,\vb{v} \cross \vb{B}_F
    \right)
\end{align}
which gives the condition on $\phi(\vb{v})$,
\begin{equation}
  \phi(-\vb{v}) \cdot \phi(\vb{v}) = 1
\end{equation}

The transformed fields approach the original ones for $\vb{v}$ going to zero,
\begin{align}
  \lim_{\vb{v} \to \vb{0}} \vb{B}_F     &= \vb{B} \\
  \lim_{\vb{v} \to \vb{0}} \phi(\vb{v}) &= 1
\end{align}
hence, for reasons of continuity, $\phi(\vb{v}) = 1$. The final form of the relativistic
electric and magnetic fields is,
\begin{align}
  \vb{B}_F &= \beta\,\vb{B}
    - (\beta-1)\ip{\vb{B}}{\tfrac{\vb{v}}{v}}\tfrac{\vb{v}}{v}
    - \frac{\beta}{c^2}\,\vb{v} \cross \vb{E} \label{eq:relativistic_maxwell} \\
  \vb{E}_F &= \beta\,\vb{E}
    - (\beta-1)\ip{\vb{E}}{\tfrac{\vb{v}}{v}}\tfrac{\vb{v}}{v}
    + \beta\,\vb{v} \cross \vb{B} \nonumber
\end{align}

\end{document}